\clearpage
\VerdierPhysicalPageAnchor{0138}%
\begingroup
\small
\setlength{\emergencystretch}{2em}%
\setlength{\parskip}{0pt}%
\setlength{\abovedisplayskip}{1pt plus 1pt minus 1pt}%
\setlength{\belowdisplayskip}{1pt plus 1pt minus 1pt}%

\begingroup\itshape
\VerdierSemanticAnchor{ch2-2-2-4-b}
\noindent b) When the category $\mathcal C$ is additive, the category
$\mathcal C(S^{-1})$ is additive and the functor $Q$ is additive.

\VerdierSemanticAnchor{ch2-2-2-4-c}
\par\noindent c) Let $\widehat S$ be the set of morphisms of $\mathcal C$
that are carried by $Q$ to isomorphisms of $\mathcal C(S^{-1})$.
\VerdierSemanticAnchor{ch2-2-2-4-c-saturation}%
The set $\widehat S$ is a saturated multiplicative system of $\mathcal C$.
Every saturated multiplicative system of $\mathcal C$ that contains $S$
contains $\widehat S$.
\endgroup

\VerdierSemanticAnchor{ch2-2-2-4-proof-b}
\par\noindent The first assertion follows immediately from the universal
property of $\mathcal C(S^{-1})$: the category
$\bigl(\mathcal C^\circ(S^{-1})\bigr)^\circ$ is a solution of this problem.
The second assertion follows immediately from assertions b) and c) of
theorem (2.2.2).
\VerdierSemanticAnchor{ch2-2-2-4-proof-c-characterization}%
It follows immediately from the description of the morphisms of
$\mathcal C(S^{-1})$ that a morphism $f$ of $\mathcal C$ is carried by $Q$
to an isomorphism if and only if there exist two morphisms $g$ and $h$ of
$\mathcal C$ such that $gf$ and $fh$ are defined and belong to $S$.
\VerdierSemanticAnchor{ch2-2-2-4-proof-c-sm2-sm3}%
It follows immediately that the set $\widehat S$ has properties SM2) and
SM3).
\VerdierSemanticAnchor{ch2-2-2-4-proof-c-sm1-sm4}%
Moreover, it follows from the definition of $\widehat S$ that this set has
properties SM1) and SM4).
\VerdierSemanticAnchor{ch2-2-2-4-proof-c-conclusion}%
Thus $\widehat S$ is a saturated multiplicative system of $\mathcal C$.
The last assertion of (2.2.4) is trivial.

\VerdierSemanticAnchor{ch2-2-2-5}
\smallskip
\noindent\textbf{2.2.5.} Let $\mathcal C$ be a category and $S$ a
multiplicative system of $\mathcal C$.
\VerdierSemanticAnchor{ch2-2-2-5-localized-category}%
The category $\mathcal C(S^{-1})$ is called the \emph{localized category}
of $\mathcal C$ with respect to the system $S$.
\VerdierSemanticAnchor{ch2-2-2-5-localization-functor}%
The functor $Q\colon\mathcal C\longrightarrow\mathcal C(S^{-1})$ is called
the \emph{localization functor}.
\VerdierSemanticAnchor{ch2-2-2-5-saturation}%
The multiplicative system $\widehat S$ is called the \emph{saturation} of
the system $S$.

\VerdierSemanticAnchor{ch2-2-2-6}
\smallskip
\noindent\textbf{Theorem 2.2.6.}\quad
\begingroup\itshape
\VerdierSemanticAnchor{ch2-2-2-6-setup}%
Let $\mathcal D$ be a triangulated category and $S$ a multiplicative system
of $\mathcal D$ (2.1.4). Denote by $\mathcal D(S^{-1})$ the localized
category of $\mathcal D$ and by
$Q\colon\mathcal D\longrightarrow\mathcal D(S^{-1})$ the localization
functor.

\VerdierSemanticAnchor{ch2-2-2-6-a}
\par\noindent a) There exists on $\mathcal D(S^{-1})$ one and only one strict
$\mathbb Z$-category structure such that $Q$ is a strict $\mathbb Z$-functor.

\VerdierSemanticAnchor{ch2-2-2-6-b}
\par\noindent b) There exists on $\mathcal D(S^{-1})$ one and only one
triangulated structure such that $Q$ is an exact functor. The distinguished
triangles for this structure are the triangles isomorphic to the images
under $Q$ of the distinguished triangles of $\mathcal D$.

\VerdierSemanticAnchor{ch2-2-2-6-c-factorization}
\par\noindent c) Every exact (respectively cohomological) functor
$F\colon\mathcal D\longrightarrow\mathcal D'$ into a triangulated
(respectively abelian) category that carries the morphisms of $S$ to
isomorphisms factors uniquely through
$Q\colon\mathcal D\longrightarrow\mathcal D(S^{-1})$ as an exact
(respectively cohomological) functor
$G\colon\mathcal D(S^{-1})\longrightarrow\mathcal D'$.
\VerdierSemanticAnchor{ch2-2-2-6-c-morphism-lifting}%
Let $F_1,F_2\colon\mathcal D\rightrightarrows\mathcal D'$ be two exact
(respectively cohomological) functors that carry the morphisms of $S$ to
isomorphisms; let
$G_1,G_2\colon\mathcal D(S^{-1})\rightrightarrows\mathcal D'$ be their
respective factorizations, and let
$m\colon F_1\longrightarrow F_2$ be a morphism of exact (respectively
cohomological) functors. There exists one and only one morphism of functors
$p\colon G_1\longrightarrow G_2$ such that $m=p\star Q$.
\endgroup

\VerdierSemanticAnchor{ch2-2-2-6-proof-a}
\par\noindent The first assertion is an immediate consequence of theorem
(2.2.2) and SM5).
\VerdierSemanticAnchor{ch2-2-2-6-proof-c-p139-boundary}%
The third assertion is an immediate consequence of the

\VerdierSemanticAnchor{ch2-2-2-6-p139-boundary}
\endgroup
